% capitulos/cap5/api.tex -- §5.10: interface HTTP (convenções e contrato dos endpoints).
% Parte do Capítulo 5 (capitulos/cap5_arquitetura.tex); os rótulos \label e os números de seção não mudaram.

\section{Interface de Programação HTTP}\label{sec:arq-api}

\subsection{Convenções}

As rotas JSON seguem um REST pragmático (Seção~\ref{sec:fund-renderizacao}). Toda rota de conteúdo e de configuração exige sessão válida, verificada no servidor antes de qualquer outra ação. Entradas e saídas são JSON em UTF-8; as rotas de escrita aceitam apenas \texttt{Content-Type: application/json}, e todas as respostas levam \texttt{Cache-Control: no-store}. Códigos de status: \texttt{200}/\texttt{201}/\texttt{204} para sucesso; \texttt{400} para JSON malformado ou \texttt{Content-Type} diferente de JSON (código \arq{VALIDATION_ERROR}); \texttt{401} sem sessão válida ou com mutação vinda de outra origem; \texttt{404} recurso inexistente; \texttt{409} conflito (por exemplo, \en{slug} duplicado); \texttt{422} falha de validação; \texttt{429} excesso de tentativas de login (resposta do better-auth); \texttt{500} erro interno. Essas convenções foram aprovadas na \adr{ADR-018}. A \adr{ADR-029} acrescenta, para o M16: \texttt{403} (recurso de outro dono ou senha provisória), \texttt{413} (corpo grande demais, recusado antes da leitura), \texttt{429} também para o limite de escritas por usuário (com \texttt{Retry-After}) e \texttt{503} (\arq{BUSY}); e o campo \texttt{reason} (1 a 500 caracteres) é obrigatório nas ações do administrador sobre outro usuário (RN19), no corpo JSON e, no \texttt{DELETE}, no parâmetro de consulta \texttt{reason}.

\textbf{Resposta de sucesso.} Toda resposta de sucesso com corpo usa o envelope \texttt{\{ "data": ... \}}, que deixa espaço para campos adicionais ao lado de \texttt{data}, como \texttt{warnings} (Seção~\ref{sec:arq-contraste}). O \texttt{POST} de criação responde \texttt{201} e o cabeçalho \texttt{Location} com o caminho do novo recurso; o \texttt{DELETE} responde \texttt{204} sem corpo; e a reordenação devolve em \texttt{data} a árvore atualizada.

\begin{lstlisting}[language=VJSON, caption={Formato padrão de sucesso}, label={lst:sucesso-json}]
{
  "data": { "id": "cm...", "title": "Introdução", "slug": "introducao", "status": "PUBLISHED" }
}
\end{lstlisting}

\textbf{Resposta de erro.} Os erros têm formato único; \texttt{details} é uma lista com um item por campo, no formato \texttt{\{ field, message \}} (validação) ou \texttt{\{ field, value \}} (valor recusado). Um \texttt{409} de \en{slug} traz em \texttt{details[0].suggestion} um \en{slug} livre entre os irmãos de destino (UC05, fluxo A1).

\begin{lstlisting}[language=VJSON, caption={Formato padrão de erro}, label={lst:erro-json}]
{
  "error": {
    "code": "SLUG_CONFLICT",
    "message": "Já existe um conteúdo com este slug na mesma coleção.",
    "details": [{ "field": "slug", "value": "introducao", "suggestion": "introducao-2" }]
  }
}
\end{lstlisting}

Códigos de erro integrados: \arq{VALIDATION_ERROR}, \arq{UNAUTHENTICATED}, \arq{NOT_FOUND}, \arq{SLUG_CONFLICT}, \arq{INVALID_PARENT}, \arq{CYCLE_DETECTED} e \arq{INTERNAL_ERROR}. Códigos novos (M16, em uso desde o merge M-013; \adr{ADR-029}): \arq{FORBIDDEN} (403), \arq{PASSWORD_CHANGE_REQUIRED} (403), \arq{QUOTA_EXCEEDED} (422), \arq{PAYLOAD_TOO_LARGE} (413), \arq{RATE_LIMITED} (429), \arq{USER_CONFLICT} (409, com \texttt{details[0].field} igual a \texttt{email} ou \texttt{handle}), \arq{LAST_ADMIN} (409), \arq{USER_HAS_CONTENT} (409) e \arq{BUSY} (503). O código \arq{VERSION_CONFLICT} (409), da RN18 (aprovada em 04/10/2026), também integra o contrato: o \texttt{PATCH} de um conteúdo recebe o campo \texttt{version} (o \texttt{updatedAt} que o autor abriu), responde 409 quando o conteúdo mudou depois, com a versão atual em \texttt{details[0].value}, e 422 quando o corpo é alterado sem o campo. Desde o merge M-032 (A-021, \adr{ADR-033}), a gravação de texto passa por um orçamento de renderização: o \texttt{POST} e o \texttt{PATCH} de conteúdo (este só quando altera o \texttt{body}, inclusive na moderação) e o \texttt{PUT} da configuração (\texttt{homeBody}) respondem 422 \arq{CONTENT_TOO_COMPLEX}, com \texttt{details[0].field} igual a \texttt{body} ou \texttt{homeBody} e a mensagem ``O texto é complexo demais para publicar: simplifique-o ou divida-o em partes.'', quando o texto não renderiza em 2~s, e nada é gravado; com a fila de validações cheia (uma em andamento e três esperando) \emph{ou} quando o mesmo usuário já tem uma validação em andamento ou na fila (no máximo uma por usuário, A-022, merge M-035), 503 \arq{BUSY} com \texttt{Retry-After: 3}. A cota em bytes (422 \arq{QUOTA_EXCEEDED}) vem \emph{antes}, porque é barata, e o limite de leitura da requisição de conteúdo é o dobro da cota mais 64~KiB (192~KiB no padrão), acima do qual a resposta é 413.

\textbf{Regras comuns de entrada.} Campos desconhecidos são descartados, inclusive \texttt{authorId}, que vem sempre da sessão (RN13). Um \texttt{PATCH} ou \texttt{PUT} sem nenhum campo conhecido responde \texttt{422}. O \en{slug} reservado também é conferido ao mover um conteúdo para a raiz (por \texttt{PATCH} ou em lote). A mutação cujo cabeçalho \texttt{Origin} indica outro site é recusada com \texttt{401} (\arq{UNAUTHENTICATED}), pois o cookie de sessão não é aceito em requisições de outra origem e o contrato não prevê \texttt{403} (Seção~\ref{sec:arq-seguranca}).

\subsubsection{Erros da Camada de Dados}\label{sec:arq-dataerror}

As funções de \arq{src/lib/db.js} lançam a classe \texttt{DataError} (com os campos \texttt{code}, \texttt{message} e \texttt{details}) quando uma regra que depende do estado do banco é violada: existência, unicidade, tipo do pai, ciclo e domínio de \texttt{type}, \texttt{status} e \texttt{order}. A validação de \emph{formato} (expressão regular e \en{slugs} reservados da RN03, URL \texttt{http}/\texttt{https} da RN17, cor \texttt{\#RRGGBB} da RN09) permanece nos validadores do Core (\arq{validators.js}). Cada Route Handler captura o \texttt{DataError} e o traduz conforme a Tabela~\ref{tab:dataerror}; qualquer outro erro vira \texttt{500} com o código \arq{INTERNAL_ERROR} e mensagem genérica, sem expor detalhes internos. As funções de escrita não chamam \arq{revalidatePath} (RN10): a revalidação é da camada de rota.

\begingroup\small
\begin{longtable}{@{}L{3.9cm}L{1.2cm}L{3.3cm}L{5.8cm}@{}}
\caption{Mapeamento dos códigos de \texttt{DataError} para a resposta HTTP}\label{tab:dataerror}\\
\toprule
\textbf{Código do \texttt{DataError}} & \textbf{HTTP} & \textbf{Código da API} & \textbf{Situação}\\
\midrule
\endfirsthead
\toprule
\textbf{Código do \texttt{DataError}} & \textbf{HTTP} & \textbf{Código da API} & \textbf{Situação}\\
\midrule
\endhead
\bottomrule
\endfoot
\arq{NOT_FOUND} & 404 & \arq{NOT_FOUND} & Conteúdo (ou, ao reordenar, algum item do lote) inexistente\\
\addlinespace
\arq{SLUG_CONFLICT} & 409 & \arq{SLUG_CONFLICT} & \en{Slug} repetido entre irmãos ou na raiz (RN03)\\
\addlinespace
\arq{PARENT_NOT_FOUND} & 422 & \arq{INVALID_PARENT} & O pai informado não existe\\
\addlinespace
\arq{PARENT_NOT_COLLECTION} & 422 & \arq{INVALID_PARENT} & O pai informado é um documento (RN02)\\
\addlinespace
\arq{CYCLE} & 422 & \arq{CYCLE_DETECTED} & O destino é o próprio nó ou um descendente (RN06)\\
\addlinespace
\arq{TYPE_HAS_CHILDREN} & 422 & \arq{VALIDATION_ERROR} & Coleção com filhos não pode virar documento (\texttt{details.field} = \texttt{type})\\
\addlinespace
\arq{AUTHOR_NOT_FOUND} & 422 & \arq{VALIDATION_ERROR} & Autor inexistente; como o autor vem da sessão, indica sessão inconsistente\\
\addlinespace
\arq{INVALID_INPUT} & 422 & \arq{VALIDATION_ERROR} & Campo obrigatório ausente ou fora do domínio (\texttt{type}, \texttt{status}, \texttt{order}, lote de reordenação)\\
\addlinespace
(outro erro) & 500 & \arq{INTERNAL_ERROR} & Falha não prevista; registrar no servidor, resposta genérica\\
\end{longtable}
\fonte{Elaborado pelos autores, a partir dos códigos de \arq{src/lib/db.js} e da sugestão do agente Dados (merge M-003).}
\endgroup

Duas observações de implementação. Em \texttt{reorderContents}, os itens do lote são aplicados em sequência e validados contra o estado resultante dos anteriores, de modo que a ordem dos itens importa quando há ciclo intermediário. E o redirecionamento de uma coleção ao primeiro documento publicado (RN14) não tem função própria: o Core usa \texttt{getContentTree} com \texttt{onlyPublished} e \texttt{resolvePath}.

\subsection{Contrato dos \emph{Endpoints}}

\begin{longtable}{@{}L{1.5cm}L{3.9cm}L{1.6cm}L{4.4cm}L{2.4cm}@{}}
\caption{Contrato das rotas HTTP (normativo)}\label{tab:api}\\
\toprule
\textbf{Método} & \textbf{Caminho} & \textbf{Sessão} & \textbf{Descrição} & \textbf{Sucesso}\\
\midrule
\endfirsthead
\toprule
\textbf{Método} & \textbf{Caminho} & \textbf{Sessão} & \textbf{Descrição} & \textbf{Sucesso}\\
\midrule
\endhead
\bottomrule
\endfoot
\texttt{GET} & \arq{/api/contents} & Sim & Árvore completa, inclusive rascunhos, sem o corpo dos documentos e com o \texttt{status} de cada nó & 200\\[1pt]
\texttt{POST} & \arq{/api/contents} & Sim & Cria conteúdo (\texttt{title}, \texttt{slug}, \texttt{type}, \texttt{parentId}, \texttt{body}, \texttt{order}, \texttt{status}; este, se omitido, assume \texttt{PUBLISHED}); o autor vem da sessão & 201 e \texttt{Location}\\[1pt]
\texttt{GET} & \arq{/api/contents/:id} & Sim & Detalhe, incluindo o corpo e o \texttt{status} & 200\\[1pt]
\texttt{PATCH} & \arq{/api/contents/:id} & Sim & Atualização parcial, inclusive de \texttt{status}; dispara revalidação (RN10, RN16) & 200\\[1pt]
\texttt{DELETE} & \arq{/api/contents/:id} & Sim & Exclui o nó e seus descendentes (RN07) & 204\\[1pt]
\texttt{PATCH} & \arq{/api/contents/reorder} & Sim & Lote \texttt{[\{id, parentId, order\}]}, em transação (RN05, RN06); devolve a árvore atualizada & 200\\[1pt]
\texttt{GET} & \arq{/api/configuration} & Sim & Configuração vigente, sem \texttt{warnings} & 200\\[1pt]
\texttt{PUT} & \arq{/api/configuration} & Sim & Atualização \emph{parcial} (campos ausentes não mudam) de identidade, cores e dos seis links sociais (RN17); logotipo e \en{favicon} aceitam URL \texttt{http(s)} ou caminho de \arq{public/} iniciado por uma única barra (D-10); devolve sempre \texttt{warnings} de contraste, calculados sobre a paleta gravada depois de gravar e revalidar (\texttt{500} se o cálculo falhar; Seção~\ref{sec:arq-contraste}) & 200\\[1pt]
\texttt{POST} & \arq{/api/auth/sign-in/email} & Não & Login (rota do better-auth), com limite de 5 tentativas por minuto por IP & 200 (429 acima do limite)\\[1pt]
\texttt{POST} & \arq{/api/auth/sign-up/email} & Não & Recusado: cadastro público desabilitado (RF22) & 400 (\arq{EMAIL_PASSWORD_SIGN_UP_DISABLED})\\[1pt]
\texttt{POST} & \arq{/api/auth/sign-out} & Sim & Encerra a sessão (RF20); o painel usa, em seu lugar, a \en{Server Function} \texttt{signOut} (\arq{admin/_lib/actions.js}) & 200\\[1pt]
\texttt{GET} & \arq{/api/auth/get-session} & Não & Sessão corrente ou nula & 200\\[1pt]
\end{longtable}
\fonte{Elaborado pelos autores. As rotas \arq{/api/auth/*} são as padrão do better-auth, conferidas na versão 1.7.6 (merge M-004); as demais rotas padrão da biblioteca (por exemplo, troca de senha) também ficam expostas em \arq{/api/auth/*} e exigiam sessão \emph{até o M16}, que as fechou: o \en{catch-all} responde só a \texttt{get-session}, \texttt{sign-in/email} e \texttt{sign-out} (Seção~\ref{sec:arq-autorizacao}).}

A leitura pública da árvore, dos documentos e da configuração \textbf{não} usa essas rotas: ocorre em Server Components, por chamadas diretas às funções de \arq{db.js}, o que mantém as páginas estáticas e aplica o filtro de publicação (RN15). As rotas de conteúdo e de configuração servem ao painel.

\textbf{Contrato do escopo multiusuário (M16 a M20; as rotas de usuários, de aparência e de conteúdo por dono estão integradas desde o merge M-013; as de notificações e a moderação, desde o merge M-015, a Parte 2 da A-014).} A Tabela~\ref{tab:api} descreve as rotas iniciais; a Tabela~\ref{tab:api-alvo} acrescenta e ajusta as do escopo multiusuário (\adr{ADR-029}), com a mesma ordem de verificação da Seção~\ref{sec:arq-autorizacao}. Sempre no envelope da \adr{ADR-018}.

\begin{longtable}{@{}L{1.5cm}L{4.1cm}L{2.1cm}L{6.1cm}@{}}
\caption{Contrato-alvo das rotas do escopo multiusuário (normativo; M16 a M20)}\label{tab:api-alvo}\\
\toprule
\textbf{Método} & \textbf{Caminho} & \textbf{Quem} & \textbf{Descrição}\\
\midrule
\endfirsthead
\toprule
\textbf{Método} & \textbf{Caminho} & \textbf{Quem} & \textbf{Descrição}\\
\midrule
\endhead
\bottomrule
\endfoot
\texttt{GET} & \arq{/api/contents} & Dono; admin com \texttt{?ownerId=} & Árvore \emph{do dono} da sessão, com rascunhos; o administrador lista a de outro usuário informando \texttt{ownerId} (moderação)\\[1pt]
\texttt{POST}, \texttt{PATCH}, \texttt{DELETE} & \arq{/api/contents}, \arq{/api/contents/:id}, \arq{/api/contents/reorder} & Dono; admin com \texttt{reason} sobre conteúdo alheio & Como na Tabela~\ref{tab:api}; o \texttt{PATCH} exige \texttt{version} e responde 409 \arq{VERSION_CONFLICT} se o conteúdo mudou (RN18); \texttt{authorId} da sessão (RN13); pai de outro dono: 422 \arq{INVALID_PARENT}, indistinguível do pai inexistente (UC05, A2); item de outro dono (o \texttt{id} da rota ou um item do lote): 403, e o lote inteiro é recusado; cotas (RN25); no conteúdo alheio, o administrador \emph{lê} desde o M16 e, desde o M-015, despublica, edita e exclui, com \texttt{reason} (1 a 500 caracteres depois de aparado: vazio ou só espaços responde 422 \arq{VALIDATION_ERROR}; no \texttt{PATCH} vai no corpo, e no \texttt{DELETE}, em \texttt{?reason=}), e o dono recebe a notificação \texttt{MODERATION} (RN26); a \texttt{version} da RN18 vale também na moderação; \texttt{parentId} e \texttt{order} no conteúdo alheio respondem 403 (mover não é moderação), e reordenar o conteúdo alheio, em \arq{/api/contents/reorder}, responde 403 mesmo ao administrador. A escrita e a notificação são uma só transação (A-017, merge M-016): se a notificação falhar, nada é gravado\\[1pt]
\texttt{GET}, \texttt{PUT} & \arq{/api/appearance} & Dono (qualquer papel) & Aparência \emph{do próprio usuário}: identidade, seis links, paleta clara e noturna; \texttt{PUT} parcial; devolve \texttt{warnings} por modo (Seção~\ref{sec:arq-contraste}); revalida só o espaço do dono. A de outro usuário responde 403 (RN19)\\[1pt]
\texttt{GET}, \texttt{PUT} & \arq{/api/configuration} & Admin & Configuração \emph{do site} (identidade, links, paletas, \texttt{homeBody}, \texttt{directoryEnabled}); usuário comum recebe 403; revalida só \texttt{/} (\adr{ADR-030})\\[1pt]
\texttt{GET}, \texttt{POST} & \arq{/api/users} & Admin & Lista, com \texttt{total}, \texttt{limit} e \texttt{offset} ao lado de \texttt{data} (\texttt{?limit=} de 1 a 100, padrão 50; \texttt{?offset=}; \texttt{?search=} por parte do e-mail); \textbf{cada usuário traz \texttt{avatarImageId}} (\texttt{null} sem avatar; \emph{alvo}, A-029), de que a interface depende para oferecer ``Remover avatar'' (RF36); cria conta (\texttt{email}, \texttt{name} de 1 a 100 caracteres, \texttt{handle}, \texttt{role}, senha opcional) com troca obrigatória e devolve a senha provisória \emph{uma única vez} em \texttt{data.temporaryPassword} (201 e \texttt{Location}; \texttt{null} se o administrador a informou)\\[1pt]
\texttt{GET}, \texttt{PATCH} & \arq{/api/users/me} & Próprio usuário & Perfil; edita o nome. Permitido com senha provisória (RN21)\\[1pt]
\texttt{POST} & \arq{/api/users/me/password} & Próprio usuário & Troca a senha (\texttt{current}, \texttt{next}; resposta \texttt{data: \{ changed: true \}}); recusa a igual à atual; revoga as outras sessões; limpa \texttt{mustChangePassword}\\[1pt]
\texttt{GET}, \texttt{PATCH}, \texttt{DELETE} & \arq{/api/users/:id} & Admin & Detalhe (com \texttt{avatarImageId}, como na lista); edita nome, \en{handle} e papel (invalida o espaço); exclui, com \texttt{?content=transfer\&to=<id>[\&onConflict=rename]} ou \texttt{?content=delete}; 409 \arq{USER_HAS_CONTENT} sem a escolha; 409 \arq{LAST_ADMIN} para o último administrador e 403 \arq{FORBIDDEN} para a própria conta (também na suspensão); \textbf{responde 200 com} \texttt{data: \{ id, content, roots, renamed \}}, exceção à convenção do 204 para que o administrador veja as raízes renomeadas (integrado no M-015)\\[1pt]
\texttt{PUT} & \arq{/api/users/:id/password} & Admin & Redefine a senha (provisória, exibida uma vez); revoga todas as sessões; notifica o usuário\\[1pt]
\texttt{POST}, \texttt{DELETE} & \arq{/api/users/:id/ban} & Admin & Suspende (\texttt{reason} obrigatório, \texttt{expiresAt} opcional) e reativa; revoga as sessões; invalida o espaço; notifica\\[1pt]
\texttt{GET} & \arq{/api/notifications} & Próprio usuário (também o administrador) & Lista as \emph{suas}, das mais recentes às mais antigas (RN26): \texttt{?unread=1} (só as não lidas e não resolvidas), \texttt{?limit=} de 1 a 100 (padrão 20) e \texttt{?cursor=} (o \texttt{nextCursor} anterior); ao lado de \texttt{data} vêm \texttt{unreadCount} (todas as não lidas, não só a página) e \texttt{nextCursor} (\texttt{null} na última). Cada item traz \texttt{title} e \texttt{text} em português, montados de \texttt{type} e \texttt{data}, sem o dono nem a chave de deduplicação. Parâmetro inválido: 422\\[1pt]
\texttt{PATCH} & \arq{/api/notifications/:id} & Próprio usuário & Exatamente uma ação: \texttt{\{ read: true \}} (idempotente) ou \texttt{\{ dismissed: true \}}; responde \texttt{data} com a ação e \texttt{unreadCount}; a de outro usuário e a inexistente respondem 404, sem distinção; não há como marcar como não lida\\[1pt]
\texttt{POST} & \arq{/api/notifications/read-all} & Próprio usuário & Marca todas as suas como lidas; responde \texttt{data: \{ count \}} e \texttt{unreadCount}\\[1pt]
\end{longtable}
\fonte{Elaborado pelos autores, a partir das ADR-023 a ADR-029 e do merge M-013. Os nomes dos parâmetros de corpo são normativos; o Core os fixa nos validadores e pode propor ajustes por \texttt{[RFC]}. Valores que o Core fixou no M-013, aceitos pelo Arquiteto: o corpo das rotas sem texto longo é limitado a 64~KiB (e o das de conteúdo, a duas vezes a cota mais 64~KiB); o administrador pode redefinir a própria senha por \texttt{PUT /api/users/:id/password}, o que revoga as próprias sessões; \texttt{homeBody} respeita \texttt{MAX\_BODY\_BYTES} (422 \arq{QUOTA_EXCEEDED}); os \texttt{details} novos são \texttt{QUOTA\_EXCEEDED} \{ \texttt{field}, \texttt{bytes} ou \texttt{total}, \texttt{max} \}, \texttt{PAYLOAD\_TOO\_LARGE} \{ \texttt{field}, \texttt{max} \} e \texttt{VERSION\_CONFLICT} \{ \texttt{field}, \texttt{value} \}. A suspensão grava a notificação \texttt{MODERATION} (\texttt{ACCOUNT\_SUSPENDED}, com o motivo e o prazo), e a reativação e a senha redefinida gravam \texttt{ACCOUNT} (\texttt{ACCOUNT\_REACTIVATED} e \texttt{PASSWORD\_RESET}); as rotas de notificações montam o texto a partir de \texttt{type} e \texttt{data} (a de contraste fica de fora até o P-013, e os eventos desconhecidos usam um texto genérico). Os \texttt{PUT} da aparência e da configuração do site devolvem, em \texttt{warnings}, os avisos de contraste dos \emph{dois} modos sobre a paleta gravada (a forma de sempre, sem chave \texttt{mode}; os da paleta noturna têm o prefixo \texttt{dark} em \texttt{field}) e geram, por modo com aviso, a notificação \texttt{CONTRAST} do dono do registro, com a chave de deduplicação \texttt{contrast:\{modo\}} na aparência e \texttt{contrast-site:\{modo\}} no site, reaberta a cada salvamento e resolvida quando o modo passa a atender (M-019). Os \texttt{PATCH} e \texttt{read-all} de notificações têm o grupo de limite próprio \texttt{notifications}, de 60 por minuto (\adr{ADR-029}). O \texttt{?cursor=} só vale se for uma notificação do próprio usuário: o de outro usuário e o inexistente (inclusive o de uma notificação dispensada entre duas páginas) respondem 422 \arq{VALIDATION_ERROR}, e o cliente recomeça pela primeira página (verificação da camada de dados, A-017).}

Os campos de entrada têm as seguintes regras de validação no servidor (\texttt{422} em caso de falha): \texttt{status} deve ser \texttt{DRAFT} ou \texttt{PUBLISHED}; cada link social, quando preenchido, deve ser uma URL absoluta com protocolo \texttt{http} ou \texttt{https}, e a cadeia vazia equivale a \texttt{null} (RN17); as cores devem ser \texttt{\#RRGGBB} (RN09); o \en{slug} segue a RN03 (1 a 80 caracteres, sem \en{slugs} reservados na raiz). A resposta de \texttt{422} traz um item de \texttt{details} por campo inválido.

\subsection{Contrato do Épico v1: Imagens, Rodapé e Aparência (Alvo)}\label{sec:arq-api-v1}

\textbf{Situação: parcialmente integrado (A-025; implementação pelas A-031, integrada em 09/10/2026, merge M-041, e A-029, o Core, ainda alvo).} As linhas \arq{/api/appearance} e \arq{/api/configuration} da Tabela~\ref{tab:api-v1} estão no código, \emph{exceto} o que depende da A-029 (a perda de \texttt{logoUrl}, \texttt{faviconUrl} e dos links sociais, que por ora ainda são aceitos e ignorados); as demais linhas são alvo. A Tabela~\ref{tab:api-v1} acrescenta e ajusta as rotas do épico v1 (\adr{ADR-034}, \adr{ADR-035}, \adr{ADR-038}); vale a mesma ordem de verificação das demais (Seção~\ref{sec:arq-autorizacao}): origem, sessão, senha provisória, papel, dono, validação. Os nomes dos parâmetros de corpo são normativos; o Core os fixa nos validadores e pode propor ajustes por \texttt{[RFC]}.

\textbf{Envio de imagem (\texttt{multipart}).} As rotas \texttt{PUT} de imagem exigem \texttt{Content-Type: multipart/form-data} com um campo \texttt{file} (os demais campos são ignorados; o nome do arquivo enviado nunca é usado); outro tipo de conteúdo responde 400 \arq{VALIDATION_ERROR}, como nas rotas JSON. O corpo é lido com limite (2~MB mais a folga fixa das fronteiras do \texttt{multipart}, que o Core determina); acima dele, 413 \arq{PAYLOAD_TOO_LARGE} \emph{antes} de decodificar. Qualquer falha de decodificação, formato fora de PNG, JPEG e WebP, arquivo vazio, dimensão acima de 25 megapixels ou tipo declarado que não confere com o conteúdo responde 422 \arq{INVALID_IMAGE} (código novo), com \texttt{details[0].field} igual a \texttt{file}, mensagem em português e \emph{sem eco} do arquivo, do nome ou do tipo enviado; nunca 500. As rotas têm limite de taxa próprio (padrão: 10 por minuto por usuário, 429 \arq{RATE_LIMITED} com \texttt{Retry-After}) e de concorrência: um processamento por vez por usuário e dois ao mesmo tempo no servidor; excedido, 503 \arq{BUSY} com \texttt{Retry-After} (\adr{ADR-034}, item 4). Uma conta com a senha provisória pendente não envia imagem (403 \arq{PASSWORD_CHANGE_REQUIRED}). O \texttt{PUT} responde 200 com \texttt{data: \{ id, url, width, height \}} (a \texttt{url} é \texttt{/media/<id>}) e o \texttt{DELETE} responde 204; remover volta ao nível de baixo da cadeia de padrões (RN30).

\begin{longtable}{@{}L{1.5cm}L{4.1cm}L{2.1cm}L{6.1cm}@{}}
\caption{Contrato-alvo das rotas do épico v1 (normativo)}\label{tab:api-v1}\\
\toprule
\textbf{Método} & \textbf{Caminho} & \textbf{Quem} & \textbf{Descrição}\\
\midrule
\endfirsthead
\toprule
\textbf{Método} & \textbf{Caminho} & \textbf{Quem} & \textbf{Descrição}\\
\midrule
\endhead
\bottomrule
\endfoot
\texttt{PUT}, \texttt{DELETE} & \arq{/api/appearance/logo}, \arq{/api/appearance/favicon} & Dono (qualquer papel) & Envia ou remove o logotipo (WebP, até 512~px no maior lado) ou o \en{favicon} (PNG, 64$\times$64) \emph{do próprio espaço}; a troca grava a imagem nova e apaga a anterior na mesma transação; revalida o espaço do dono. A de outro usuário responde 403 (RN19)\\[1pt]
\texttt{PUT}, \texttt{DELETE} & \arq{/api/configuration/logo}, \arq{/api/configuration/favicon} & Admin & O mesmo, para o \emph{sistema}; revalida o portal e os espaços que herdam a imagem (\adr{ADR-038}, item 6); usuário comum recebe 403\\[1pt]
\texttt{PUT}, \texttt{DELETE} & \arq{/api/users/me/avatar} & Próprio usuário & Envia ou remove o avatar (WebP, 256$\times$256, corte quadrado centralizado); revalida o espaço do dono e o portal (o avatar entra no diretório)\\[1pt]
\texttt{DELETE} & \arq{/api/users/:id/avatar} & Admin & Remove o avatar de outro usuário; \texttt{reason} obrigatório (1 a 500 caracteres); apaga a imagem e grava a notificação \texttt{MODERATION} (\texttt{data.event} igual a \texttt{AVATAR\_REMOVED}, com o motivo) na \emph{mesma transação}; 404 se o usuário não tem avatar; 204\\[1pt]
\texttt{GET}, \texttt{PUT} & \arq{/api/footer-links} & Dono (qualquer papel) & Rodapé \emph{do próprio espaço}. \texttt{PUT} com \texttt{\{ links: [\{ icon, url \}] \}} substitui a lista inteira, em transação, e a ordem do vetor é a ordem do rodapé; de 0 a 8 itens; \texttt{icon} da lista fechada (RN17), sem repetir; \texttt{url} \texttt{https://} ou, no ícone \texttt{email}, \texttt{mailto:} com e-mail válido; um item inválido recusa o lote com 422 \arq{VALIDATION_ERROR} e um \texttt{details} por item (\texttt{field} como \texttt{links[2].url}); \texttt{GET} responde \texttt{data: [\{ icon, url \}]} na ordem; revalida o espaço do dono\\[1pt]
\texttt{GET}, \texttt{PUT} & \arq{/api/configuration/footer-links} & Admin & O mesmo, para o rodapé do \emph{sistema} (o portal, o \texttt{/login} e o painel do administrador); revalida o portal\\[1pt]
\texttt{GET} & \arq{/media/:id} & Público & Serve a imagem. \emph{Fora de} \texttt{/api}, sem sessão e sem \texttt{no-store}: \texttt{Content-Type} fixo da saída, \texttt{X-Content-Type-Options: nosniff}, \texttt{ETag} (o \texttt{sha256}) e \texttt{Cache-Control: public, max-age=31536000, immutable}; \texttt{If-None-Match} responde 304; identificador inexistente responde 404, sem distinguir. Serve também a imagem de um espaço suspenso (RN23, \adr{ADR-034}, item 8)\\[1pt]
\texttt{GET}, \texttt{PUT} & \arq{/api/appearance} & Dono (qualquer papel) & \emph{Integrado (M-041), exceto o último item.} Como na Tabela~\ref{tab:api-alvo}, com os campos novos: \texttt{showLogo} e \texttt{showName} (ao menos um verdadeiro, senão 422 com \texttt{field} igual a \texttt{showName}) e a \emph{herança}: \texttt{siteName} e as seis cores aceitam \texttt{null} (herdam o sistema; o nome vazio ou só de espaços vira \texttt{null}); a paleta de um modo tem as três cores definidas ou as três nulas (senão 422 com \texttt{field} igual a \texttt{lightPalette} ou \texttt{darkPalette}). O \texttt{GET} devolve a aparência \emph{crua} (com os nulos) em \texttt{data} e, ao lado, \texttt{effective} (a mescla sobre o sistema, sem o padrão embutido, e sem a chave \texttt{inherits}) e \texttt{inherits} (\texttt{lightPalette}, \texttt{darkPalette}, \texttt{siteName}, \texttt{logo}, \texttt{favicon}). O \texttt{PUT} responde \texttt{\{ data, warnings \}} e os avisos e a notificação de contraste valem só para o modo cuja paleta o \emph{próprio usuário} definiu (\adr{ADR-038}, item 7): quem herda uma paleta ruim do sistema não é avisado, e voltar a herdar um modo resolve a notificação dele. \emph{Alvo (A-029):} \texttt{logoUrl}, \texttt{faviconUrl} e os links sociais deixam de existir (por ora ainda são aceitos e ignorados)\\[1pt]
\texttt{GET}, \texttt{PUT} & \arq{/api/configuration} & Admin & \emph{Integrado (M-041), exceto o logotipo e o \en{favicon}.} Como na Tabela~\ref{tab:api-alvo}, com os campos novos: \texttt{showLogo}, \texttt{showName} (a mesma regra), \texttt{modeIconPair} (um dos quatro pares da RN32) e \texttt{directoryTitle} (1 a 40 pontos de código, texto simples); \emph{alvo (A-029):} perde \texttt{logoUrl}, \texttt{faviconUrl} e os links sociais (por ora ainda aceitos e ignorados). As paletas e o nome continuam obrigatórios (\texttt{null} responde 422), porque o sistema é a base da herança. Mudar a paleta de um modo ou o nome do sistema revalida também os espaços que herdam o que mudou, e mudar o \texttt{modeIconPair} revalida o espaço de todos (\adr{ADR-038}, item 6; Seção~\ref{sec:arq-isr}); o logotipo e o \en{favicon} do sistema entram com a A-029. Os avisos de contraste são do administrador, dono da paleta do sistema\\[1pt]
\end{longtable}
\fonte{Elaborado pelos autores, a partir das ADR-034, ADR-035 e ADR-038 e do documento-fonte do épico (\S4). O \texttt{GET /api/notifications?limit=8} da Tabela~\ref{tab:api-alvo} alimenta o painel suspenso do sino (RF40) e não muda.}

A \texttt{url} que o servidor guarda no rodapé é só texto: a interface a renderiza como \texttt{href} depois de conferir o esquema, com \texttt{target="\_blank"} e \texttt{rel="noopener noreferrer"} nos \texttt{https} (RN17). O corpo de \texttt{PUT /api/footer-links} respeita o limite comum de 64~KiB. As rotas de envio de imagem têm cabeçalhos de resposta como as demais (\texttt{Cache-Control: no-store}); só \texttt{GET /media/:id} tem o cabeçalho imutável.
